%% Lecture 11 -- Field transformations and the electrostatic limit
\section{Lecture 11 -- Field Transformations and Electrostatics}\label{sec:lec11}

\begin{center}
\begin{tabular}{ll}
  \textbf{Date:}\quad 12/05/2026 & \\
\end{tabular}
\end{center}
\hrule\vspace{1.5em}

\subsection*{Overview}

This lecture closes one conceptual loop from Lecture~10 and opens the next topic.
The central message is that the electric and magnetic fields are not two independent
relativistic objects.  They are different components of one antisymmetric tensor
$F_{\mu\nu}$, so a Lorentz boost generally mixes them.

The lecture has four parts:
\begin{itemize}
  \item a clarification of the invariant $E^2-B^2$: it is compared at the same
        spacetime event, not at unrelated spatial points;
  \item a proof that the four-dimensional Levi-Civita symbol $\epsilon_{\mu\nu\rho\sigma}$
        behaves as a tensor under proper Lorentz transformations, and as a pseudotensor
        if reflections are included;
  \item a derivation of the Lorentz transformation laws for $\bE$ and $\bB$;
  \item the beginning of electrostatics, defined as the time-reversal symmetric sector
        of Maxwell theory.
\end{itemize}

The physical bridge is simple: after Maxwell's equations have been written in covariant
form, the next useful question is how familiar non-covariant language, such as
``electric field'' and ``magnetic field'', is seen by different inertial observers.

%------------------------------------------------------------
\subsection{What Does an Electromagnetic Invariant Compare?}\label{subsec:lec11-invariant-events}

\subsubsection*{Physical Motivation}
An invariant is not a statement that a field has the same value everywhere in space.
Fields can vary from point to point.  The invariant statement compares what two observers
assign to the \emph{same event} in spacetime.  This distinction matters because a Lorentz
transformation changes the coordinates of the event, not the event itself.

Recall from Lecture~10 that the scalar
\begin{equation}\label{eq:lec11:E2minusB2}
  \boxed{
  \frac{1}{2}F_{\mu\nu}F^{\mu\nu}=B^2-E^2
  }
\end{equation}
is Lorentz invariant, equivalently $E^2-B^2$ is invariant up to the conventional
overall sign used in writing the contraction.
\textit{Physical significance:} If two observers describe the same spacetime event
$x^\mu=(ct,\br)$ and $x'^\mu=(ct',\br')$, they may disagree about the separate
values of $\bE$ and $\bB$, but they agree on the scalar combination $E^2-B^2$ at
that event.

\begin{center}
\begin{tikzpicture}[scale=1.0, >=Stealth]
  \draw[->] (-0.2,0) -- (5.0,0) node[right] {$x$};
  \draw[->] (0,-0.2) -- (0,3.1) node[above] {$ct$};
  \draw[->, blue!70!black, thick] (0.4,0.15) -- (3.8,2.7) node[right] {$ct'$};
  \draw[->, blue!70!black, thick] (0.1,0.45) -- (4.5,1.4) node[right] {$x'$};
  \fill[red!70!black] (2.6,1.55) circle (2.2pt);
  \node[red!70!black, right] at (2.6,1.55) {same event};
  \draw[dashed] (2.6,0) node[below] {$x$} -- (2.6,1.55) -- (0,1.55) node[left] {$ct$};
  \node[align=center] at (2.8,-0.65) {Different coordinates,\\same physical point in spacetime.};
\end{tikzpicture}
\end{center}

%------------------------------------------------------------
\subsection{The Levi-Civita Symbol as a Tensor and Pseudotensor}\label{subsec:lec11-epsilon}

\subsubsection*{Physical Motivation}
In three dimensions, cross products and oriented volumes rely on the Levi-Civita symbol.
In four-dimensional electrodynamics the analogous object $\epsilon_{\mu\nu\rho\sigma}$
is needed to build the second electromagnetic invariant and the dual field tensor.
Before using it in relativistic equations, we must know how it transforms.

Define the completely antisymmetric object by
\begin{equation}\label{eq:lec11-epsilon-def}
  \boxed{
  \epsilon_{0123}=+1,\qquad
  \epsilon_{\mu\nu\rho\sigma}=0
  \quad\text{if any two indices are equal,}
  }
\end{equation}
with all other nonzero components fixed by antisymmetry.
\textit{Physical significance:} $\epsilon_{\mu\nu\rho\sigma}$ is the algebraic
device that remembers orientation in spacetime, just as the determinant remembers
the oriented volume scaling of a matrix.

To test whether it is a tensor, apply the tensor transformation rule:
\begin{equation}\label{eq:lec11-epsilon-transform-question}
  X_{\mu_1\mu_2\mu_3\mu_4}
  =
  \Lambda^{\nu_1}{}_{\mu_1}
  \Lambda^{\nu_2}{}_{\mu_2}
  \Lambda^{\nu_3}{}_{\mu_3}
  \Lambda^{\nu_4}{}_{\mu_4}
  \epsilon_{\nu_1\nu_2\nu_3\nu_4}.
\end{equation}
The object $X_{\mu_1\mu_2\mu_3\mu_4}$ is completely antisymmetric.  Indeed, if
two free indices are interchanged, the corresponding two dummy indices in
$\epsilon_{\nu_1\nu_2\nu_3\nu_4}$ can be interchanged, and the antisymmetry of
$\epsilon$ gives a minus sign.  In four dimensions, every completely antisymmetric
rank-four object is proportional to $\epsilon_{\mu_1\mu_2\mu_3\mu_4}$, so
\begin{equation}\label{eq:lec11-epsilon-alpha}
  X_{\mu_1\mu_2\mu_3\mu_4}
  =
  \alpha\,
  \epsilon_{\mu_1\mu_2\mu_3\mu_4}.
\end{equation}
To find $\alpha$, choose $(\mu_1,\mu_2,\mu_3,\mu_4)=(0,1,2,3)$:
\begin{align}
  \alpha
  &=
  \Lambda^{\nu_1}{}_{0}
  \Lambda^{\nu_2}{}_{1}
  \Lambda^{\nu_3}{}_{2}
  \Lambda^{\nu_4}{}_{3}
  \epsilon_{\nu_1\nu_2\nu_3\nu_4}
  \notag\\
  &=
  \det \Lambda .
  \label{eq:lec11-alpha-det}
\end{align}
Therefore
\begin{equation}\label{eq:lec11-epsilon-transform}
  \boxed{
  \Lambda^{\nu_1}{}_{\mu_1}
  \Lambda^{\nu_2}{}_{\mu_2}
  \Lambda^{\nu_3}{}_{\mu_3}
  \Lambda^{\nu_4}{}_{\mu_4}
  \epsilon_{\nu_1\nu_2\nu_3\nu_4}
  =
  (\det\Lambda)\,
  \epsilon_{\mu_1\mu_2\mu_3\mu_4}.
  }
\end{equation}
\textit{Physical significance:} For proper Lorentz transformations, $\det\Lambda=+1$,
so $\epsilon_{\mu\nu\rho\sigma}$ behaves as a tensor.  Under a reflection, such as
parity, $\det\Lambda=-1$, so it gains an extra sign.  That is why it is often called
a pseudotensor when improper transformations are allowed.

This also explains why angular momentum in relativistic notation is naturally an
antisymmetric rank-two tensor rather than an ordinary four-vector: in more than
three dimensions, a rotation is best associated with a plane, not only with an axis.

%------------------------------------------------------------
\subsection{Lorentz Boost Setup}\label{subsec:lec11-boost-setup}

\subsubsection*{Physical Motivation}
The cleanest way to see the mixing of $\bE$ and $\bB$ is to compare two inertial
frames related by a boost along the $x$-axis.  This loses no conceptual content:
any boost can be decomposed into a rotation, a boost along one axis, and another
rotation.

Let $K'$ move with speed $v$ along the positive $x$ direction as seen from $K$.
Write
\begin{equation}\label{eq:lec11-beta-gamma}
  \boxed{
  \beta=\frac{v}{c},\qquad
  \gamma=\frac{1}{\sqrt{1-\beta^2}} .
  }
\end{equation}
\textit{Physical significance:} $\beta$ measures the boost speed in units of the
maximum signal speed $c$, while $\gamma$ is the relativistic stretching factor that
controls time dilation, length contraction, and field mixing.

The coordinate transformation from primed coordinates to unprimed coordinates is
\begin{align}
  x &= \gamma(x'+v t'), &
  ct &= \gamma(ct'+\beta x'), &
  y&=y', &
  z&=z'.
  \label{eq:lec11-lorentz-coordinates}
\end{align}
Equivalently, for the four-potential
\begin{equation}\label{eq:lec11-four-potential}
  A^\mu=(\phi,\Avec),
\end{equation}
the same boost gives
\begin{align}
  \phi &= \gamma(\phi' + \beta A'_x),&
  A_x &= \gamma(A'_x+\beta\phi'),&
  A_y &= A'_y,&
  A_z &= A'_z .
  \label{eq:lec11-potential-transform}
\end{align}

The fields are defined by
\begin{equation}\label{eq:lec11-fields-potentials}
  \boxed{
  \bE = -\nabla\phi-\frac{1}{c}\pdv{\Avec}{t},
  \qquad
  \bB = \nabla\times\Avec .
  }
\end{equation}
\textit{Physical significance:} The potentials transform simply as a four-vector,
but the observed fields involve spacetime derivatives of the potentials.  This is
why the transformation of $\bE$ and $\bB$ contains both the boost and the curl-like
structure of electromagnetism.

%------------------------------------------------------------
\subsection{Transforming the Electric Field}\label{subsec:lec11-electric-transform}

\subsubsection*{Physical Motivation}
We first transform $\bE$ directly from the potentials.  This derivation is algebraic,
but its physical content is important: a purely magnetic part of the field tensor
in one frame can contribute to the electric field measured in another frame.

The $y$ component in $K$ is
\begin{equation}\label{eq:lec11-Ey-start}
  E_y
  =
  -\pdv{\phi}{y}
  -\frac{1}{c}\pdv{A_y}{t}.
\end{equation}
Since $y=y'$, the first derivative is
\begin{equation}\label{eq:lec11-y-derivative}
  \pdv{}{y}=\pdv{}{y'}.
\end{equation}
For the time derivative, the inverse boost gives
\begin{equation}\label{eq:lec11-inverse-time-chain}
  \pdv{}{t}
  =
  \pdv{t'}{t}\pdv{}{t'}
  +
  \pdv{x'}{t}\pdv{}{x'}
  =
  \gamma\pdv{}{t'}
  -\gamma v\pdv{}{x'} .
\end{equation}
Substituting \eqref{eq:lec11-potential-transform} into \eqref{eq:lec11-Ey-start},
\begin{align}
  E_y
  &=
  -\pdv{}{y'}\!\left[\gamma(\phi'+\beta A'_x)\right]
  -\frac{1}{c}
  \left(\gamma\pdv{}{t'}-\gamma v\pdv{}{x'}\right)A'_y
  \notag\\
  &=
  \gamma\left(-\pdv{\phi'}{y'}-\frac{1}{c}\pdv{A'_y}{t'}\right)
  +\gamma\beta\left(\pdv{A'_y}{x'}-\pdv{A'_x}{y'}\right)
  \notag\\
  &=
  \gamma\left(E'_y+\beta B'_z\right).
  \label{eq:lec11-Ey-result}
\end{align}
In the last line we used
\begin{equation}\label{eq:lec11-Bz-def}
  B'_z=(\nabla'\times\Avec')_z
  =
  \pdv{A'_y}{x'}-\pdv{A'_x}{y'} .
\end{equation}
Similarly,
\begin{align}
  E_x &= E'_x, \label{eq:lec11-Ex-result}\\
  E_z &= \gamma(E'_z-\beta B'_y). \label{eq:lec11-Ez-result}
\end{align}
Thus, for a boost along the $x$ direction,
\begin{equation}\label{eq:lec11-E-components}
  \boxed{
  E_x=E'_x,\qquad
  E_y=\gamma(E'_y+\beta B'_z),\qquad
  E_z=\gamma(E'_z-\beta B'_y).
  }
\end{equation}
\textit{Physical significance:} The component parallel to the boost is unchanged,
while the perpendicular electric field mixes with the perpendicular magnetic field.

It is useful to write this in vector form.  With the sign convention of
\eqref{eq:lec11-lorentz-coordinates}, let $\boldsymbol{\beta}_{K|K'}$ denote the
velocity of $K$ as seen from $K'$, divided by $c$.  Then
\begin{equation}\label{eq:lec11-E-vector}
  \boxed{
  \bE_{\parallel}=\bE'_{\parallel},\qquad
  \bE_{\perp}
  =
  \gamma\left(\bE'_{\perp}
  +\boldsymbol{\beta}_{K|K'}\times\bB'\right).
  }
\end{equation}
\textit{Physical significance:} The cross product says that only fields transverse
to the relative motion are mixed; a field already pointing along the boost direction
has no transverse area element with which to mix.

%------------------------------------------------------------
\subsection{Transforming the Magnetic Field}\label{subsec:lec11-magnetic-transform}

\subsubsection*{Physical Motivation}
The magnetic field can be transformed by repeating the potential calculation, but
the tensor method shows the deeper reason immediately.  The magnetic components
are spatial components of $F^{\mu\nu}$, and a boost mixes spatial and temporal
indices.  Since temporal-spatial components are electric-field components, $\bB$
must mix with $\bE$.

For the electromagnetic tensor,
\begin{equation}\label{eq:lec11-F-transform}
  F^{\mu\nu}
  =
  \Lambda^\mu{}_{\alpha}\Lambda^\nu{}_{\beta}F'^{\alpha\beta}.
\end{equation}
Using the convention from Lecture~10,
\begin{equation}\label{eq:lec11-F-components}
  F^{0i}=-E_i,\qquad
  F^{i0}=E_i,\qquad
  F^{ij}=-\epsilon^{ijk}B_k .
\end{equation}
Take the component $F^{12}=-B_z$.  Since the boost only mixes the $0$ and $1$
indices, while the $2$ direction is untouched,
\begin{align}
  F^{12}
  &=
  \Lambda^1{}_{\alpha}\Lambda^2{}_{\beta}F'^{\alpha\beta}
  \notag\\
  &=
  \Lambda^1{}_{0}F'^{02}+\Lambda^1{}_{1}F'^{12}
  \notag\\
  &=
  \gamma\beta(-E'_y)+\gamma(-B'_z).
  \label{eq:lec11-F12-work}
\end{align}
Therefore
\begin{equation}\label{eq:lec11-Bz-result}
  B_z=\gamma(B'_z+\beta E'_y).
\end{equation}
The other components are obtained in the same way:
\begin{equation}\label{eq:lec11-B-components}
  \boxed{
  B_x=B'_x,\qquad
  B_y=\gamma(B'_y-\beta E'_z),\qquad
  B_z=\gamma(B'_z+\beta E'_y).
  }
\end{equation}
\textit{Physical significance:} The magnetic field transforms in the same pattern
as the electric field, but with the opposite sign in the mixing term.  This is the
tensor statement that electric and magnetic fields are two faces of one object.

In vector form,
\begin{equation}\label{eq:lec11-B-vector}
  \boxed{
  \bB_{\parallel}=\bB'_{\parallel},\qquad
  \bB_{\perp}
  =
  \gamma\left(\bB'_{\perp}
  -\boldsymbol{\beta}_{K|K'}\times\bE'\right).
  }
\end{equation}
\textit{Physical significance:} If an observer moves relative to a purely electric
configuration, they generally measure a magnetic field.  Magnetism is therefore
not an optional extra phenomenon; it is required by relativity once electric fields
exist.

The inverse transformations are obtained by replacing $\boldsymbol{\beta}_{K|K'}$
with $-\boldsymbol{\beta}_{K|K'}$.  The component formulas are often the safest
way to check signs.

%------------------------------------------------------------
\subsection{Why Use Both Languages?}\label{subsec:lec11-which-language}

\subsubsection*{Physical Motivation}
Tensor notation and three-vector notation answer different practical questions.
Tensor notation exposes symmetry; three-vector notation is often better for solving
ordinary boundary-value problems.

The tensor $F_{\mu\nu}$ makes it clear that:
\begin{equation}\label{eq:lec11-one-field}
  \boxed{
  \text{there is one electromagnetic field, not two unrelated fields.}
  }
\end{equation}
\textit{Physical significance:} What one observer calls ``electric'' and ``magnetic''
depends on their state of motion.  The invariant object is the tensor, not the split
into $\bE$ and $\bB$.

However, in electrostatics there is no need to carry all tensor components if only
$\bE$ is present in the chosen frame.  In that setting, the scalar potential $\phi$
and the equation for $\bE$ are the efficient language.

%------------------------------------------------------------
\subsection{The Covariant Form of the Homogeneous Maxwell Equations}\label{subsec:lec11-homogeneous}

\subsubsection*{Physical Motivation}
Lecture~10 wrote the inhomogeneous Maxwell equations as
$\partial_\mu F^{\mu\nu}=(4\pi/c)J^\nu$.  The other two Maxwell equations also have
a covariant form.  They do not come from sources; they come from the definition of
$F_{\mu\nu}$ in terms of the potential.

The homogeneous Maxwell equations are
\begin{equation}\label{eq:lec11-homogeneous-covariant}
  \boxed{
  \partial_\alpha F_{\beta\gamma}
  +
  \partial_\beta F_{\gamma\alpha}
  +
  \partial_\gamma F_{\alpha\beta}
  =
  0 .
  }
\end{equation}
\textit{Physical significance:} This equation is the covariant form of
$\nabla\cdot\bB=0$ and $\nabla\times\bE=-(1/c)\partial_t\bB$.  It says that the
field tensor is built from a potential, so the antisymmetrized second derivative
vanishes.

Indeed, with
\begin{equation}\label{eq:lec11-F-from-A}
  F_{\mu\nu}=\partial_\mu A_\nu-\partial_\nu A_\mu,
\end{equation}
substitution into \eqref{eq:lec11-homogeneous-covariant} gives pairs of second
derivatives such as $\partial_\alpha\partial_\beta A_\gamma$ and
$-\partial_\beta\partial_\alpha A_\gamma$.  Since partial derivatives commute,
each pair cancels:
\begin{equation}\label{eq:lec11-second-derivative-cancel}
  \partial_\alpha\partial_\beta A_\gamma
  -
  \partial_\beta\partial_\alpha A_\gamma
  =
  0 .
\end{equation}

%------------------------------------------------------------
\subsection{Electrostatics as a Symmetry Sector}\label{subsec:lec11-electrostatics}

\subsubsection*{Physical Motivation}
Electrostatics is often introduced as ``nothing depends on time''.  That is true,
but it is not the most structural definition.  The lecture defines electrostatics
as the sector of Maxwell theory that respects time reversal $t\to -t$.  This
distinguishes it from magnetostatics, where currents can be steady but reverse
under time reversal.

The Maxwell equations in Gaussian units are
\begin{align}
  \nabla\cdot\bB &= 0, \label{eq:lec11-maxwell-divB}\\
  \nabla\times\bE &= -\frac{1}{c}\pdv{\bB}{t}, \label{eq:lec11-maxwell-faraday}\\
  \nabla\cdot\bE &= 4\pi\rho, \label{eq:lec11-maxwell-gauss}\\
  \nabla\times\bB &= \frac{4\pi}{c}\Jvec+\frac{1}{c}\pdv{\bE}{t}. \label{eq:lec11-maxwell-ampere}
\end{align}
The sources must also obey
\begin{equation}\label{eq:lec11-continuity}
  \boxed{
  \pdv{\rho}{t}+\nabla\cdot\Jvec=0 .
  }
\end{equation}
\textit{Physical significance:} Charge conservation is the consistency condition
that allows the sources to appear on the right-hand side of Maxwell's equations.

Under time reversal, velocities reverse, so currents reverse:
\begin{equation}\label{eq:lec11-time-reversal-current}
  \Jvec \longrightarrow -\Jvec .
\end{equation}
The electric field is even under time reversal, while the magnetic field is odd:
\begin{equation}\label{eq:lec11-time-reversal-fields}
  \bE\longrightarrow \bE,\qquad
  \bB\longrightarrow -\bB .
\end{equation}
If a solution is invariant under time reversal for any choice of the time origin,
then there can be no genuine time dependence, no current, and no magnetic field:
\begin{equation}\label{eq:lec11-electrostatic-conditions}
  \boxed{
  \pdv{}{t}=0,\qquad
  \Jvec=0,\qquad
  \bB=0 .
  }
\end{equation}
\textit{Physical significance:} Electrostatics is not merely ``slow dynamics''.
It is the time-reversal-even part of the theory: charges sit as prescribed sources,
and the magnetic, current-carrying sector is absent.

With \eqref{eq:lec11-electrostatic-conditions}, Maxwell's equations reduce to
\begin{equation}\label{eq:lec11-electrostatic-maxwell}
  \boxed{
  \nabla\times\bE=0,\qquad
  \nabla\cdot\bE=4\pi\rho .
  }
\end{equation}
\textit{Physical significance:} The first equation says the electric field is
conservative; the second says charge density is the source of electric flux.

Since $\nabla\times\bE=0$, one can write
\begin{equation}\label{eq:lec11-E-potential}
  \boxed{
  \bE=-\nabla\phi .
  }
\end{equation}
\textit{Physical significance:} In electrostatics the field can be described by
one scalar function.  Equipotential surfaces organize the problem, and the field
points in the direction of steepest decrease of $\phi$.

Substituting \eqref{eq:lec11-E-potential} into Gauss's law gives Poisson's equation:
\begin{equation}\label{eq:lec11-poisson}
  \boxed{
  \nabla^2\phi=-4\pi\rho .
  }
\end{equation}
\textit{Physical significance:} Electrostatics becomes a boundary-value problem:
given charges and boundary conditions, solve one elliptic equation for the scalar
potential.

\begin{center}
\begin{tikzpicture}[scale=1.0, >=Stealth]
  \fill[red!70!black] (0,0) circle (3pt) node[below left] {$+$};
  \foreach \a in {0,30,...,330} {
    \draw[->, red!70!black] (0,0) -- ({2.0*cos(\a)},{2.0*sin(\a)});
  }
  \draw[blue!70!black, dashed] (0,0) circle (0.75);
  \draw[blue!70!black, dashed] (0,0) circle (1.35);
  \draw[blue!70!black, dashed] (0,0) circle (1.95);
  \node[blue!70!black] at (2.55,1.0) {$\phi=\text{const.}$};
  \node[red!70!black] at (2.75,-0.55) {$\bE=-\nabla\phi$};
\end{tikzpicture}
\end{center}

The lecture stops at the reduction to the electrostatic sector.  The next step is
to solve \eqref{eq:lec11-poisson} for concrete charge distributions and boundary
conditions.

%------------------------------------------------------------
\subsection{Summary}\label{subsec:lec11-summary}

\begin{itemize}
  \item Electromagnetic invariants compare fields at the same spacetime event as
        described by different inertial observers.
  \item The four-dimensional Levi-Civita symbol satisfies
        $\epsilon\mapsto(\det\Lambda)\epsilon$; it is a tensor for proper Lorentz
        transformations and a pseudotensor if reflections are included.
  \item A boost mixes perpendicular electric and magnetic fields:
        $\bE_\perp$ receives a contribution from $\bB'$, and $\bB_\perp$ receives
        a contribution from $\bE'$.
  \item The tensor $F_{\mu\nu}$ is the invariant electromagnetic object; the split
        into $\bE$ and $\bB$ is observer-dependent.
  \item Electrostatics is the time-reversal-even sector of Maxwell theory, where
        $\Jvec=0$, $\bB=0$, and $\bE=-\nabla\phi$ with
        $\nabla^2\phi=-4\pi\rho$.
\end{itemize}
